\documentclass[leqno, a4paper, 12pt]{jsarticle}
\usepackage{amssymb}
\usepackage{amsthm}
\usepackage{amsmath}
\usepackage{comment}
\usepackage{url}

\usepackage{enumitem}
%\usepackage{showkeys}


%定理環境 thmはあまり使わずpropをなるべく使う。
%主張は基本的に証明の中で使い、そのため番号を付けない。
%注意環境を付けてみた。
\newtheorem{thm}{定理}
\newtheorem{prop}[thm]{命題}
\newtheorem{cor}[thm]{系}
\newtheorem{lem}[thm]{補題}
\newtheorem{df}[thm]{定義}
\newtheorem{exam}[thm]{例}
\newtheorem{fact}[thm]{事実}
\newtheorem*{claim}{主張}
\newtheorem*{cau}{注意}
\newtheorem*{rmk}{注意}
\renewcommand{\proofname}{証明}
%矢印たち。
%\toは\rightarrowと同じ。\getsは\leftarrowと同じ。同値は\iff
\newcommand{\To}{\Longrightarrow}
\newcommand{\Gets}{\LongLeftarrow}
%黒板太字
\newcommand{\nn}{\mathbb{N}}
\newcommand{\zz}{\mathbb{Z}}
\newcommand{\qq}{\mathbb{Q}}
\newcommand{\rr}{\mathbb{R}}
\newcommand{\cc}{\mathbb{C}}
%無限大
\newcommand{\mgn}{\infty}
%差集合は\setminusで統一する。等号の否定は\neq
%真部分集合は\subsetneqq　偏微分は\partial
%ディスプレイスタイル。\dfracでディスプレイ分数
\newcommand{\dpst}{\displaystyle}

\title{論文メモ
\large{\\--- 同相写像の拡張 ---}
}
\author{ナンブ　キトラ}
\date{\today}
\begin{document}
\maketitle

本棚を整理したいので，
溜まっている論文のメモを書いて未来への手紙とすることにする．

この論文の束は
多分同相写像の拡張の話に関するものだと思う．


論文
\cite{antoine1924voisinages}
は大切なことをしているような気がするが
フランス語なのでよくわからなかった．

論文
\cite{lavrentieff1924contribution}
もフランス語なのでよくわからず．

論文
\cite{MR0029505}
はおそらく代数トポロジー的な量を使って
射影平面の間のペアノ連続体の間の同相写像が拡張できるための必要十分条件を決定している．

論文
\cite{MR0004767}
では空間が有限次元あるいは無限次元の平行体のときに閉集合上の同相写像が拡張できるための十分条件を与えている．


論文
\cite{knaster1953notion}
はフランス語で書かれているのでよくわからず．
噂に(\cite{MR0320992})
よるとカントール集合のコンパクト集合同士の同相写像が
全体に拡張できることを証明しているようだ．

論文
\cite{MR0320992}
ではコンパクト部分集合同士の
同相写像が全体に拡張できるための十分条件を求めている．
応用としてヒルベルト空間$\ell^{2}$の有理点全体の空間
$R_{\omega}$, つまり時として エルデシュ空間と呼ばれる１次元空間だが，この空間のコンパクト部分集合同士の同相は全体に拡張できるらしい．


論文
\cite{MR0930078}
はコンパクト部分集合同士の
同相写像を拡張する話である．
これはなんかすごいことをやっている気がするが全然わからなかった．ごめん．




%READ!
%myplain is the same to amsplain style. 
\nocite{*}
\bibliographystyle{myplaindoidoi}
\bibliography{homeoe.bib}



\end{document}